% Propietario conservado: /Users/ruben/Documents/New project/output/EXTREMA_MEDIA_RAZON_RECIPROCIDAD_ES_EN_20260918/ARTICULO/ES/source/libro/02_norma_ternaria.tex
\chapter{Ternary algebraic norm and optimal recovery radii}
\label{x_reciprocidad:lib02:norma-ternaria}

The recovery radius of the two memories contains the factor \(\sqrt{73}\). Its origin can be traced without introducing a target magnitude: it comes from the Gram matrix of the nonadic reader and, in another realization of the ternary cycle already constructed, is the scale factor of a similitude. Its square, \(73\), is the index of the rank-two lattice inclusion. These occurrences are brought together through explicit operations; they are not identified through a decimal resemblance to alpha. The subsequent comparison with \(10^4\alpha_A\) will be treated as an inverse problem, with its data, permitted transformations, and exact conclusions.

The raw material preserves its APP--TRIT--TPK production and the subsequent
cut described in \cref{x_reciprocidad:lib01:registro}. The integer register, the
orientation of shifts, the charge and the readers are received
before distances are calculated. The canonical case uses \(8+1\) and the
denominator nine. The parameter \(a>0\) introduced below
is a subsequent algebraic family, not a new seed or a physical
parameter selected by alpha. Errors belong to the reading channel;
they are not identified with admissible events or histories.

\section{Parametric double-memory operators and Gram matrix}
\label{x_reciprocidad:lib02:familia-memoria}

In \(V=\mathbb R^{12}\), let \((Sx)_i=x_{i+1}\), with indices modulo twelve. For \(a>0\) and \(j=3,4\), define
\begin{equation}
 T_{j,a}=\frac{aI+S^j}{a+1},\qquad
 D_{j,a}=\frac{\sqrt a(S^j-I)}{a+1},\qquad
 M_ax=(D_{3,a}x,D_{4,a}T_{3,a}x).
 \label{x_reciprocidad:lib02:lectores}
\end{equation}
The second memory is produced after the first compression. Archiving both is not equivalent to reinjecting the first complement into a reversible two-channel evolution. Let
\[
 p(a)=a^2+a+1,\quad A(a)=a(a^2+1),\quad B(a)=a^2,
 \qquad\Pi=I-\frac{\mathbf1\mathbf1^*}{12},\quad P_u=I-\Pi.
\]

\begin{proposition}[Balance and difference form]
\label{x_reciprocidad:lib02:gram}
For every \(a>0\),
\[
 T_{j,a}^*T_{j,a}+D_{j,a}^*D_{j,a}=I,\qquad
 G_a:=M_a^*M_a=I-(T_{4,a}T_{3,a})^*(T_{4,a}T_{3,a}),
\]
and
\begin{align}
 (a+1)^4G_a={}&4ap(a)I
 -A(a)(S^3+S^4+S^8+S^9)\notag\\
 &-B(a)(S+S^5+S^7+S^{11}).
 \label{x_reciprocidad:lib02:gram-explicito}
\end{align}
Equivalently,
\begin{align}
 (a+1)^4\|M_ah\|^2={}&A(a)\sum_i
 \bigl((h_i-h_{i+3})^2+(h_i-h_{i+4})^2\bigr)\notag\\
 &+B(a)\sum_i
 \bigl((h_i-h_{i+1})^2+(h_i-h_{i+5})^2\bigr).
 \label{x_reciprocidad:lib02:energia-aristas}
\end{align}
\end{proposition}
\begin{proof}
Orthogonality of \(S\) gives
\[
 T_{j,a}^*T_{j,a}
 =\frac{(a^2+1)I+a(S^j+S^{-j})}{(a+1)^2},\qquad
 D_{j,a}^*D_{j,a}
 =\frac{a(2I-S^j-S^{-j})}{(a+1)^2}.
\]
Their sum is \(I\); applying the identity to the two stages cancels the intermediate term \(T_{3,a}^*T_{3,a}\). Since all operators are polynomials in \(S\), the product of the two individual Gram matrices has diagonal \((a^2+1)^2\), single terms of weight \(a(a^2+1)\), and cross terms of weight \(a^2\). The identity \((a+1)^4-(a^2+1)^2=4a(a^2+a+1)\) gives \eqref{x_reciprocidad:lib02:gram-explicito}. Expanding the squared differences gives \eqref{x_reciprocidad:lib02:energia-aristas}; each unoriented edge of each step appears exactly once, and the weighted degree is \(4A+4B=4ap(a)\).
\end{proof}

At \(a=8\) we obtain
\begin{align}
 6561G_8={}&2336I-520(S^3+S^4+S^8+S^9)\notag\\
 &-64(S+S^5+S^7+S^{11}),
 \label{x_reciprocidad:lib02:gram-nonadico}
\end{align}
whose first row is \((2336,-64,0,-520,-520,-64,0,-64,-520,-520,0,-64)/6561\). Also,
\[
 T_{4,8}T_{3,8}=\frac{64I+8S^3+8S^4+S^7}{81}.
\]
The four distinct shifts produce the diagonal contribution
\[
 1-\frac{64^2+8^2+8^2+1}{81^2}
 =\frac{2336}{6561}
 =\frac{32(8^2+8+1)}{9^4}.
\]

\begin{theorem}[Kernel, spectrum, and complete inverse]
\label{x_reciprocidad:lib02:inversa-memorias}
For every \(a>0\), \(\ker M_a=\mathbb R\mathbf1\). The eigenvalues of \(G_a\), indexed by the twelfth-root modes \(m\), are
\[
\begin{array}{c|c|c}
 m&\lambda_m(G_a)&\text{multiplicity}\\\hline
 0&0&1\\
 3,9&2a/(a+1)^2&2\\
 4,8&3a/(a+1)^2&2\\
 6&4a/(a+1)^2&1\\
 1,5,7,11&a(5a^2+4a+5)/(a+1)^4&4\\
 2,10&a(7a^2+2a+7)/(a+1)^4&2
\end{array}
\]
In particular,
\begin{equation}
 L_a^{\rm rec}=(G_a+P_u)^{-1}M_a^*,\qquad
 L_a^{\rm rec}M_a=\Pi,\qquad
 \|L_a^{\rm rec}\|=\frac{a+1}{\sqrt{2a}}.
 \label{x_reciprocidad:lib02:inversa}
\end{equation}
\end{theorem}
\begin{proof}
If \(M_ah=0\), the first memory implies \(S^3h=h\), hence \(T_{3,a}h=h\); the second implies \(S^4h=h\). Since \(S=S^4(S^3)^{-1}\), the vector is uniform. The converse is immediate. In mode \(m\), the balance gives
\[
 \lambda_m(G_a)=1-
 \frac{[a^2+1+2a\cos(\pi m/2)]
 [a^2+1+2a\cos(2\pi m/3)]}{(a+1)^4}.
\]
The twelve cosine pairs give the table. Relative to \(2a/(a+1)^2\), the differences of the last two values are
\[
 \frac{3a(a^2+1)}{(a+1)^4}>0,\qquad
 \frac{a(5a^2-2a+5)}{(a+1)^4}>0.
\]
The other positive values are larger multiples. Hence the minimum on \(\mathbf1^\perp\) is \(2a/(a+1)^2\). The operator \(G_a+P_u\) equals one on the uniform direction and is positive and invertible on its complement. The formula displayed is the least-squares inverse, and its norm is the reciprocal of the smallest nonzero singular value of \(M_a\).
\end{proof}

At \(a=1\), \(T_{3,1}=(I+S^3)/2\) is singular, but \eqref{x_reciprocidad:lib02:inversa} remains valid. Singularity of an individual compression does not eliminate recovery through its two memories. At \(a=8\), the norm of the inverse is \(9/4\).

\section{Minimum distance in the domain of integer registers}
\label{x_reciprocidad:lib02:redes}

The algebraic decoding domains are
\[
 \mathcal C_0(a)=M_a\mathbb Z^{12},\qquad
 \Lambda_q=\{k\in\mathbb Z^{12}:\mathbf1^*k=q\},\qquad
 \mathcal C_q(a)=M_a\Lambda_q.
\]
The first image is identified with the lattice \(\Pi\mathbb Z^{12}\)
or with classes modulo \(\mathbb Z\mathbf1\); the second preserves a
fixed integer charge. These are domains containing the integer publications
of the indicated type. Their use does not assert that each candidate corresponds to a
produced HMT history.

\begin{lemma}[Minimum cut of the memory graph]
\label{x_reciprocidad:lib02:cortes}
Every nonempty proper cut of the twelve positions has at least four edges of steps \(1,5\) and at least four of steps \(3,4\). Its weight in \eqref{x_reciprocidad:lib02:energia-aristas} is therefore at least \(4A(a)+4B(a)=4ap(a)\); a cut consisting of one position attains this bound.
\end{lemma}
\begin{proof}
Each step \(1\) and \(5\) forms a connected cycle of twelve positions;
each cycle crosses a nontrivial cut at least twice. For the steps
\(3,4\), if the set is invariant under neither, each family
contributes at least two edges. If it is invariant under \(S^3\), it is a union
of classes modulo three; each of the four triangles of \(S^4\)
crosses it twice, giving eight edges. If it is invariant under \(S^4\),
it is a union of classes modulo four; each of the three length-four cycles
of \(S^3\) crosses it at least twice, giving at least six.
Simultaneous invariance would force invariance under \(S\), impossible
for a nonempty proper set. All weights are positive; the
cut of a single position has exactly four edges of each
weight.
\end{proof}

\begin{theorem}[Exact distances and fixed-charge minimizers]
\label{x_reciprocidad:lib02:distancias}
For every \(a>0\) and every integer charge \(q\),
\begin{equation}
 d_0(a)^2=\min_{h\notin\mathbb Z\mathbf1}\|M_ah\|^2
 =\frac{4ap(a)}{(a+1)^4},\qquad
 d_q(a)^2=\frac{8ap(a)}{(a+1)^4}.
 \label{x_reciprocidad:lib02:distancias-exactas}
\end{equation}
In the first minimum, \(h\) ranges over \(\mathbb Z^{12}\); the second is the minimum distance between distinct words of \(\mathcal C_q(a)\). The differences attaining the second minimum are exactly \(e_i-e_j\) with \(j-i\equiv2,6,10\pmod{12}\).
\end{theorem}
\begin{proof}
For a nonconstant integer vector \(h\), \((h_i-h_j)^2\geq|h_i-h_j|\). Each absolute difference decomposes as a sum of cuts of the level sets \(\{i:h_i>t\}\). There is at least one nontrivial level; by \cref{x_reciprocidad:lib02:cortes}, the weighted energy is at least \(4ap(a)\). The vector \(e_i\) attains this energy. Division by \((a+1)^4\) proves the first minimum over all integer vectors, not only over indicators of subsets.

For fixed charge, the difference \(h\ne0\) is integral and has zero sum. Its squared norm is even because \(h_i^2\equiv h_i\pmod2\). If \(\|h\|^2\geq4\), the spectral minimum gives
\[
 \|M_ah\|^2\geq\frac{8a}{(a+1)^2}
 >\frac{8ap(a)}{(a+1)^4},
\]
since \((a+1)^2-p(a)=a>0\). Only \(\|h\|^2=2\) remains, that is, \(h=e_i-e_j\). Its energy is \(2((G_a)_{ii}-(G_a)_{ij})\). The off-diagonal entries are zero at separations \(2,6,10\), negative with weight \(B(a)/(a+1)^4\) at \(1,5,7,11\), and negative with weight \(A(a)/(a+1)^4\) at \(3,4,8,9\). The diagonal is \(4ap(a)/(a+1)^4\); this proves the minimum and the exhaustive classification over the entire half-line \(a>0\).
\end{proof}

\begin{theorem}[Strict correction radii of the memories]
\label{x_reciprocidad:lib02:radios-memoria}
For a datum \(y=M_ak+e\), with error measured in the joint norm of the two memories, the optimal uniform radii are
\begin{equation}
 r_0(a)=\frac{\sqrt{ap(a)}}{(a+1)^2},\qquad
 r_q(a)=\frac{\sqrt{2ap(a)}}{(a+1)^2}.
 \label{x_reciprocidad:lib02:radios}
\end{equation}
If \(\|e\|<r_0(a)\), the nearest neighbour in \(\mathcal C_0(a)\) recovers \(\Pi k\) exactly. With known charge \(q\) and \(\|e\|<r_q(a)\), the nearest neighbour in \(\mathcal C_q(a)\) recovers \(k\) exactly.
\end{theorem}
\begin{proof}
The lattice \(\Pi\mathbb Z^{12}\subset(1/12)\mathbb Z^{12}\) is discrete, and \(M_a\) is bounded below on \(\mathbf1^\perp\). Its images are discrete and admit a nearest neighbour. If \(c'\) is another codeword and \(d\) is the corresponding minimum distance, \(\|y-c'\|\geq d-\|e\|>\|e\|\) when \(\|e\|<d/2\). The minimum is unique, and the preimage is unique in each indicated domain. The midpoint of two words at minimum distance lies at \(d/2\) from both: no decoder can uniformly distinguish both inputs from that same datum. Optimality refers to the complete algebraic domains, not to an additional subset of histories that might have greater separation.
\end{proof}

In the canonical case,
\begin{equation}
 r_0(8)=\frac{2\sqrt{146}}{81}\simeq0.298346814162829,
 \qquad
 r_q(8)=\frac{4\sqrt{73}}{81}\simeq0.421926110879878.
 \label{x_reciprocidad:lib02:radio73}
\end{equation}
The \(4\) comes from \(\sqrt{2\cdot8}\) when taking half the distance between fixed-charge registers; \(73=8^2+8+1\) comes from the Gram matrix; \(81=9^2\) comes from the two nonadic normalizations. The same product \(9\times9\) counts the APP cells, but these two roles are not identified without their typed map. Nor does the factor four alone prove an assertion about spatial directions.

If the memories are stored as \(z=M_ak/\sqrt a\), the noise norm changes. At \(a=8\), the radii for this rational datum are \(r_0/\sqrt8=\sqrt{73}/81\) and \(r_q/\sqrt8=\sqrt{146}/81\). Using the radii in \eqref{x_reciprocidad:lib02:radio73} without this conversion changes the problem. The application to \(Q_K\), \(\eta_K\), \(B_K\), and the incidence screen has been proved in \cref{x_reciprocidad:lib01:correccion-incidencia}; in particular, the exact threshold \(K_4=729\) does not prevent its discrete restoration.

\section{Finite decoder and parameter duality}
\label{x_reciprocidad:lib02:computacion}

\begin{proposition}[Floor and ceiling without a target table]
\label{x_reciprocidad:lib02:decodificador}
With known charge, below radius \(r_q(a)\), it suffices to examine at most \(2^{12}=4096\) integer candidates. Without charge, below \(r_0(a)\), it suffices to examine at most \(12\cdot2^{12}=49152\) candidates distributed among twelve charge residues. The residual minimum returns \(k\) or its uniform class, respectively.
\end{proposition}
\begin{proof}
With charge \(q\), form \(\widetilde k=L_a^{\rm rec}y+(q/12)\mathbf1\). Then
\[
 \|\widetilde k-k\|
 <\|L_a^{\rm rec}\|r_q(a)
 =\frac{\sqrt{p(a)}}{a+1}<1.
\]
Each true integer coordinate belongs to \(\{\lfloor\widetilde k_i\rfloor,
\lceil\widetilde k_i\rceil\}\). These products are enumerated, those whose sum differs from \(q\) are removed, and \(\|M_ak'-y\|\) is minimized. The correct candidate is present, and \cref{x_reciprocidad:lib02:radios-memoria} ensures its uniqueness. Coordinates that are already integers reduce the number of candidates.

Without charge, write \(\rho\in\{0,\ldots,11\}\) for the charge residue. Each class modulo \(\mathbb Z\mathbf1\) has a unique representative whose sum is \(\rho\), obtained by subtracting \(\lfloor q/12\rfloor\mathbf1\). For each \(\rho\), form \(\widetilde k_\rho=L_a^{\rm rec}y+(\rho/12)\mathbf1\). At the correct residue,
\[
 \|\widetilde k_\rho-k_\rho\|
 <\|L_a^{\rm rec}\|r_0(a)
 =\frac{\sqrt{p(a)}}{\sqrt2(a+1)}<1.
\]
Apply floor and ceiling, retain the sum \(\rho\), and minimize the residual across the twelve groups. Uniqueness of the code class gives a unique centred output. These procedures use the reader and integrality; they do not receive \(K\) as a comparison table.
\end{proof}

For rational \(a\), rational arithmetic can be retained by writing \(R_a=M_a/\sqrt a\), \(z=y/\sqrt a\). The centred reconstruction is
\[
 (R_a^*R_a+P_u)^{-1}R_a^*z=L_a^{\rm rec}y.
\]
The residual can be compared through its rational squares. These guarantees concern twelve coordinates: they assert neither polynomial complexity for arbitrary dimension, cryptographic authentication, nor quantum advantage. The perturbed example in \cref{x_reciprocidad:lib01:correccion-incidencia} illustrates the canonical algorithm.

\begin{proposition}[Reciprocal duality and maximum radius]
\label{x_reciprocidad:lib02:dualidad-a}
The substitution \(a\mapsto1/a\) preserves the Gram matrix, distances, radii, and norm of the inverse. We have
\[
 D_{j,1/a}=D_{j,a},\qquad
 T_{j,1/a}=S^jT_{j,a}^*.
\]
The maximum radii of this family are attained only at \(a=1\):
\[
 r_q(a)\leq\frac{\sqrt6}{4},\qquad
 r_0(a)\leq\frac{\sqrt3}{4}.
\]
\end{proposition}
\begin{proof}
The two operator identities follow by multiplying numerator and denominator by \(a\). The individual Gram matrices and formula \eqref{x_reciprocidad:lib02:gram} are invariant under reciprocity. Put \(t=a/(a+1)^2\in(0,1/4]\). Then
\[
 r_q(a)^2=2t(1-t),\qquad r_0(a)^2=t(1-t).
\]
Both expressions increase for \(0<t\leq1/4\), and \(t=1/4\) is equivalent to \(a=1\). Reciprocity leaves \(t\) fixed. As \(a\) tends to zero or infinity, \(t\) and the radii tend to zero.
\end{proof}

Equality of Gram matrices does not mean that \(M_a=M_{1/a}\): the second
complement includes a different compression. The endpoints \(a=0\)
and \(a=\infty\) do not belong to the domain of the theorem. Optimizing the
family at \(a=1\), or using the reciprocal representative \(1/8\),
does not retroactively change the canonical origin \(a=8\).

\section{The norm on the ternary augmentation module}
\label{x_reciprocidad:lib02:modulo-ternario}

The cycle \(r\mapsto4r\pmod9\), with orbits \((1,4,7)\) and
\((2,8,5)\), provides the augmentation module whose chart and form are
\[
 C=\begin{pmatrix}0&-1\\1&-1\end{pmatrix},\qquad
 G_{A_2}=\begin{pmatrix}2&-1\\-1&2\end{pmatrix},\qquad
 C^2+C+I=0,\qquad C^TG_{A_2}C=G_{A_2}.
\]
The action and its contragredient orientation belong to the antecedent ternary fibres. On this already constructed module, we compose
\[
 Q_-(a)=aI-C,\qquad Q_+(a)=(a+1)I+C.
\]

\begin{proposition}[Ternary similitude and integer index]
\label{x_reciprocidad:lib02:norma-A2}
For \(a>0\),
\[
 Q_-Q_+=p(a)I,\qquad Q_-^{\dagger_{G_{A_2}}}=Q_+,
 \qquad Q_-^TG_{A_2}Q_-=p(a)G_{A_2},\qquad \det Q_-=p(a).
\]
At \(a=8\), both conjugate charts are
\[
 Q_-=
 \begin{pmatrix}8&1\\-1&9\end{pmatrix},\qquad
 Q_+=\begin{pmatrix}9&-1\\1&8\end{pmatrix},\qquad Q_-Q_+=73I.
\]
Each integer quotient is cyclic of order 73. On the quotient of \(Q_-\), \(C\) acts by \(8\); on that of \(Q_+\), by \(-9\equiv64\equiv8^{-1}\pmod{73}\).
\end{proposition}
\begin{proof}
Expanding the product and substituting \(C^2=-C-I\) gives \(p(a)I\). Invariance of the form implies \(C^{\dagger_{G_{A_2}}}=C^{-1}=C^2=-I-C\), from which the adjoint follows. Composing product and adjoint gives the similitude. The determinant of \(\bigl(\begin{smallmatrix}a&1\\-1&a+1\end{smallmatrix}\bigr)\) is \(p(a)\). At \(a=8\), the greatest common divisor of the entries of each matrix is one and its determinant is 73; the Smith normal form is \(\operatorname{diag}(1,73)\). The relation \(8I-C=0\) in the first quotient and \(9I+C=0\) in the second gives the indicated actions. Since \(8^3\equiv1\pmod{73}\) and \(8\not\equiv1\), both actions have order three and opposite orientations.
\end{proof}

Thus \(\sqrt{73}\) is a similitude scale factor in the inherited ternary metric, not merely a number inside a radius. Its subsequent complex recognition is \(|8-\omega|^2=73\), with \(\omega^2+\omega+1=0\). The complex unit describes the realization of the operator; it was not used to select the register. These equalities do not identify the operator \(Q_-\) with \(K\), with the entire TPK dynamics, or with the cokernel of the action-scale reader.

\section{Invariant contribution of the complete cycle}
\label{x_reciprocidad:lib02:ciclo-completo}

On the twelve positions, let \(U=S^4\). There are now four cycles of length three. Their projectors are
\[
 P_0=\frac{I+U+U^2}{3},\qquad P=I-P_0.
\]
The rank of \(P_0\) is four and that of \(P\) is eight. In particular, \(P_0\ne P_u\), since the latter preserves only the global uniform direction; nor is \(P\) to be confused with the exceptional projector \(P_3\). We shall now use \(Q_{-,a}=aI-U\) and \(Q_{+,a}=(a+1)I+U\) for the operators of the complete cycle.

\begin{theorem}[Ternary Gram matrix with a fixed sector]
\label{x_reciprocidad:lib02:gram-ternario}
We have
\begin{equation}
 Q_{-,a}^*Q_{-,a}
 =(a-1)^2P_0+p(a)P=p(a)I-3aP_0.
 \label{x_reciprocidad:lib02:gram-Qmenos}
\end{equation}
For \(v\ne0\), let
\[
 w(v)=\frac{\|P_0v\|^2}{\|v\|^2},\qquad
 s_{-,a}(v)=\frac{\|Q_{-,a}v\|^2}{\|v\|^2}.
\]
Then \(s_{-,a}(v)=p(a)-3aw(v)\). In the canonical case,
\begin{equation}
 Q_-^*Q_-=49P_0+73P=73I-24P_0,
 \qquad s(v)=73-24w(v)\in[49,73].
 \label{x_reciprocidad:lib02:73-24w}
\end{equation}
\end{theorem}
\begin{proof}
From \(U^3=I\) and \(U^*=U^2\) it follows that \(P_0\) is the orthogonal projector onto \(\ker(U-I)\). Expanding the Gram matrix,
\[
 Q_{-,a}^*Q_{-,a}=(a^2+1)I-a(U+U^2).
\]
The relation \(U+U^2=3P_0-I\) gives \eqref{x_reciprocidad:lib02:gram-Qmenos}. Evaluating on \(v\) and dividing by its squared norm proves the remaining identities.
\end{proof}

The term \(24\) is exactly \(3\cdot8\): it is the difference \(p(8)-(8-1)^2\) between the augmentation reading and the fixed reading. It already appears in a single cycle of three positions. For \(m\) cycles, the fixed rank is \(m\), the dimension is \(3m\), and
\[
 \operatorname{tr}(Q_{-,a}^*Q_{-,a})=3m(a^2+1).
\]
At \(m=4,a=8\), the trace is 780 and the total trace correction is \(24\cdot4=96\). The numerical equality \(24=2\cdot12\) is not the operator origin of the coefficient.

The two complete charts also distinguish their determinants. In a
single cycle,
\[
 \det(8I-U)=8^3-1=511=7\cdot73,
 \qquad
 \det(9I+U)=9^3+1=730=10\cdot73.
\]
The fixed factor is 7 and 10, respectively; the augmentation module contributes 73. On the twelve positions, \(\det(8I-U)=511^4=68184176641\). The identity
\begin{equation}
 73=8^2+8+1=9^2-9+1=\frac{9^3+1}{9+1}
 =\frac{729+1}{10}.
 \label{x_reciprocidad:lib02:completacion73}
\end{equation}
is preserved. Its interpretation as an index and determinant factor precedes any subsequent completion or refinement reading.

\section{Six received states and conservation of the complete reading}
\label{x_reciprocidad:lib02:seis-estados}

We again fix \(a=8\) and omit that subscript in the readers. The six objects being compared are
\[
 K,\quad k=\Pi K,\quad m_0=D_3K,\quad m_1=D_4T_3K,
 \quad u=P_3k,\quad q_K=Q_KK=k-u.
\]
None is selected from the target of the test. To give the evaluations without replacing them by decimals, write \(a_v=\|v\|^2\), \(b_v=\|P_0v\|^2\). We obtain
\begin{equation}
\begin{array}{c|c|c}
 v&a_v&b_v\\\hline
 K&4251257&10684363/3\\
 k&11789915/12&1170761/4\\
 m_0&15413392/81&16838032/243\\
 m_1&1105759232/6561&0\\
 u&1327717/4+(568896/5)\sqrt5
   &3029657/20+(145120/3)\sqrt5\\
 q_K&1951691/3-(568896/5)\sqrt5
   &2292404/15-(145120/3)\sqrt5
\end{array}
 \label{x_reciprocidad:lib02:tabla-normas}
\end{equation}
The table is obtained by applying \(P_0=(I+S^4+S^8)/3\), the formulas for \(D_3,D_4T_3\), and the finite character projector \eqref{x_reciprocidad:lib01:P3} to \eqref{x_reciprocidad:lib01:K}. All its entries belong to \(\mathbb Q(\sqrt5)\); the first four are rational. In particular, the first three weights and readings are
\begin{align*}
 w(K)&=\frac{10684363}{12753771},&
 s(K)&=\frac{224866857}{4251257},\\
 w(k)&=\frac{3512283}{11789915},&
 s(k)&=\frac{776369003}{11789915},\\
 w(m_0)&=\frac{1052377}{2890011},&
 s(m_0)&=\frac{61904585}{963337}.
\end{align*}
For \(u\) and \(q_K\), the exact expressions are \(w=b_v/a_v\), \(s=73-24b_v/a_v\), with the numerators and denominators in \eqref{x_reciprocidad:lib02:tabla-normas}. The subsequent evaluations are
\[
\begin{array}{c|c|c}
 v&w(v)\text{, approximately}&s(v)\text{, approximately}\\\hline
 K&0.8377414805&52.89420446705527\\
 k&0.2979057101&65.85026295779062\\
 m_0&0.3641429046&64.26057028848679\\
 m_1&0&73\text{ exact}\\
 u&0.4428244530591254&62.37221312658099\\
 q_K&0.1127385160275117&70.29427561533972
\end{array}
\]
The equality for the second memory is forced: \(P_0D_4T_3=\sqrt8P_0(U-I)T_3/9=0\), while \(\|m_1\|^2>0\). It is not an evaluational coincidence.

\begin{proposition}[Complete memory does not alter the fixed weight]
\label{x_reciprocidad:lib02:memoria-ternaria}
The map
\[
 V_2v=(t(v),m_0(v),m_1(v))
 =(T_4T_3v,D_3v,D_4T_3v)
\]
is isometric and intertwines \(P_0\) and \(Q_-=8I-U\) with their diagonal copies. It separately preserves \(\|v\|^2\), \(\|P_0v\|^2\), and \(\|Q_-v\|^2\). Hence the joint readings \(w\) and \(s\) do not change.
\end{proposition}
\begin{proof}
The two-stage balance gives \(V_2^*V_2=I\). Each channel is a polynomial in \(S\), so it commutes with \(U,P_0,Q_-\). Applying the isometry to \(v\), \(P_0v\), and \(Q_-v\) gives the three norm equalities. After division, the total reading is the average of the channel readings weighted by their squared norms, not their unweighted arithmetic mean. Zero channels do not contribute to the total ratio.
\end{proof}

The terminal component explicitly confirms this partition. For \(K\),
\[
 \|t(K)\|^2=\frac{25538253193}{6561},\qquad
 \|P_0t(K)\|^2=\frac{848595371}{243},
\]
\[
 w(t(K))=\frac{22912075017}{25538253193},\qquad
 s(t(K))=\frac{1314402682681}{25538253193}
 \simeq51.46799480558349.
\]
The sum of the three squared norms is \(4251257\), and that of the fixed-sector norms is \((848595371+16838032)/243=10684363/3\). For \(k=\Pi K\),
\[
 \|t(k)\|^2=\frac{16367568169}{26244},\qquad
 \|P_0t(k)\|^2=\frac{217142795}{972}.
\]
The complements do not change under centring, and their sums recover the norms in the second row of \eqref{x_reciprocidad:lib02:tabla-normas}. Discarding the terminal component or a complement alters the partial reading; retaining them introduces no hidden correction.

In a complete infinite memory, the same conclusion holds when the channels commute with \(U\): the positive limit and its square root also commute with it. Without this commutation, the transport \(\mathcal VP_0\mathcal V^*\) is used on the image, in accordance with \cref{x_reciprocidad:lib01:isometria}, not an assumed diagonal copy.

\section{Inversion, charts and bilateral balance}
\label{x_reciprocidad:lib02:bidireccional}

\begin{proposition}[What inversion or a frame change preserves]
\label{x_reciprocidad:lib02:inversion-carta}
Replacing \(U\) by \(U^{-1}\) exchanges each \(Q_{\pm,a}\) with its adjoint and preserves \(P_0\), their Gram matrices, singular values, determinants, and induced distances. For every \(v\ne0\), it preserves \(w(v)\) and \(s_{-,a}(v)\). This is not the same as inverting \(Q_{-,a}\), whose inverse, for \(a\ne1\), is
\begin{equation}
 Q_{-,a}^{-1}
 =\frac1{a-1}P_0+\frac1{p(a)}Q_{+,a}P.
 \label{x_reciprocidad:lib02:inversa-Q}
\end{equation}
For an invertible chart \(R\), transporting
\(U'=RUR^{-1}\), \(v'=Rv\) and
\(g'=(R^{-1})^*R^{-1}\) preserves these norm ratios in
the metric \(g'\).
\end{proposition}
\begin{proof}
Since \(U^*=U^{-1}=U^2\), inversion produces the adjoints.
The operators are normal polynomials in \(U\); their Gram matrices depend
on \(U+U^2\), which does not change. The written inverse is verified by
sectors: \(Q_{-,a}\) acts by \(a-1\) on the fixed sector and the product
with \(Q_{+,a}\) equals \(p(a)I\) on the augmentation. For the chart,
\(Q'v'=RQv\) and \(R^*g'R=I\); numerator and denominator recover
exactly their original values. The adjoint with respect to \(g'\)
is also the transport of the original adjoint.
\end{proof}

Reversing the order of the twelve positions satisfies \(RSR^{-1}=S^{-1}\) and preserves \(P_0\). The translations \(S^j\), including the antipodal translation \(S^6\), commute with \(U\); a global negation also leaves the norms unchanged. None of these operations alone turns 73 into a different reading.

With the Euclidean metric fixed, an orthogonal frame applied only to the state preserves \(s_{-,a}(v)\) for all \(v\) if and only if \(R^*P_0R=P_0\), by \eqref{x_reciprocidad:lib02:gram-Qmenos} and polarization. The group of these transformations is \(O(\operatorname{im}P_0)\times O(\operatorname{im}P)
\cong O(4)\times O(8)\), larger than the commutant of \(U\). By contrast, a state filter \(c_0P_0+c_1P\), when its output is nonzero, produces
\begin{equation}
 w'=
 \frac{|c_0|^2w}{|c_0|^2w+|c_1|^2(1-w)}.
 \label{x_reciprocidad:lib02:filtro-pesos}
\end{equation}
The reading may change. Choosing these coefficients from a target is
a declared adjustment of the filter, not a chart change of the same
reading.

\begin{theorem}[Bilateral balance and explicit recovery]
\label{x_reciprocidad:lib02:balance-bilateral}
For every \(a>0\), the complete cycle satisfies
\begin{align}
 Q_{+,a}-Q_{-,a}^*&=3P_0,&
 Q_{+,a}Q_{-,a}&=p(a)I-3P_0,\notag\\
 Q_{+,a}^*Q_{+,a}&=p(a)I+3(a+1)P_0,
 \label{x_reciprocidad:lib02:dos-cartas}
\end{align}
and
\begin{equation}
 (a+1)Q_{-,a}^*Q_{-,a}
 +aQ_{+,a}^*Q_{+,a}=(2a+1)p(a)I.
 \label{x_reciprocidad:lib02:balance-pesado}
\end{equation}
The map
\begin{equation}
 W_av=\frac{(\sqrt{a+1}\,Q_{-,a}v,
                    \sqrt a\,Q_{+,a}v)}{\sqrt{(2a+1)p(a)}}
 \label{x_reciprocidad:lib02:W}
\end{equation}
is isometric and is inverted on its image by \(W_a^*\), of norm one. For the unnormalized channels \(x=Q_{-,a}v\), \(y=Q_{+,a}v\),
\begin{equation}
 v=\frac{x+y}{2a+1},\qquad
 Uv=\frac{ay-(a+1)x}{2a+1}.
 \label{x_reciprocidad:lib02:reconstruccion-bilateral}
\end{equation}
A pair \((x,y)\) belongs to this unnormalized image if and only if \(Q_{+,a}x=Q_{-,a}y\).
\end{theorem}
\begin{proof}
The expression \(Q_+-Q_-^*\) contains \(I+U+U^2=3P_0\). Expanding \(Q_+Q_-\), \(Q_+^*Q_+\), and using \(U^2=3P_0-I-U\) gives \eqref{x_reciprocidad:lib02:dos-cartas}. The weighted combination cancels the \(P_0\) terms in \eqref{x_reciprocidad:lib02:gram-Qmenos} and \eqref{x_reciprocidad:lib02:dos-cartas}; this proves \(W_a^*W_a=I\). Since \(Q_-+Q_+=(2a+1)I\), the sum of the channels gives \(v\); their combination \(ay-(a+1)x\) gives \((2a+1)Uv\).

The image condition is necessary because the two operators commute. If it holds, define \(v=(x+y)/(2a+1)\). Then
\[
 Q_-v=\frac{Q_-x+Q_-y}{2a+1}
 =\frac{(Q_-+Q_+)x}{2a+1}=x,
\]
and similarly \(Q_+v=y\). The adjoint inverse formula is that of any isometry; for the raw pair, the linear inverse \((x,y)\mapsto(x+y)/(2a+1)\) has norm \(\sqrt2/(2a+1)\) on the product space.
\end{proof}

At \(a=8\), the fixed and augmentation sectors have readings \((49,100)\) and \((73,73)\), respectively, and
\[
 9\|Q_-v\|^2+8\|Q_+v\|^2=17\cdot73\|v\|^2=1241\|v\|^2,
 \qquad v=\frac{Q_-v+Q_+v}{17}.
\]
For \(s_{+,a}(v)=\|Q_{+,a}v\|^2/\|v\|^2\),
\[
 s_{+,a}=p(a)+3(a+1)w,
 \qquad (a+1)s_{-,a}+a s_{+,a}=(2a+1)p(a).
\]
Cancellation holds even with a fixed component. The bilateral composition is a formalization on antecedent operators, not an assertion that every TPK return is automatically this pair.

If \(\mathcal V:H\to\mathcal M\) is a complete isometric memory, \(U^{\mathcal V}=\mathcal VU\mathcal V^*\) satisfies \((U^{\mathcal V})^3=\Pi_{\mathcal M}\), where \(\Pi_{\mathcal M}=\mathcal V\mathcal V^*\) is the identity on the image. There the transported Gram matrix is
\[
 (Q_{-,a}^{\mathcal V})^*Q_{-,a}^{\mathcal V}
 =p(a)\Pi_{\mathcal M}-3a\mathcal VP_0\mathcal V^*.
\]
It preserves the readings and bilateral recovery. Using diagonal copies of \(U\) is legitimate only when the channels intertwine that action; otherwise transport by \(\mathcal V\) is the correct definition on the image.

\section{Fine-structure coordinate and integer iterations}
\label{x_reciprocidad:lib02:ensayo-inverso}

The complete alpha reader, with carry, supplies the block prefix
\[
 (007,297,352,569,283,800,997,285,105,472,380,662,999,\ldots).
\]
The zero-boundary window ending in 663 is a different object; it is not interchanged with the complete output. The inherited closed interval suffices for the test:
\begin{equation}
 \frac{7297352569283800997285105472380662}{10^{36}}
 \leq\alpha_A\leq
 \frac{7297352569283800997285105472380663}{10^{36}}.
 \label{x_reciprocidad:lib02:intervalo-alpha}
\end{equation}
The constant is not identified with either of its rational endpoints. Put \(s_\alpha=10^4\alpha_A\), so that
\begin{equation}
 \begin{gathered}
 \ell\leq s_\alpha\leq h,\\
 \ell:=72.97352569283800997285105472380662,\\
 h:=72.97352569283800997285105472380663.
 \end{gathered}
 \label{x_reciprocidad:lib02:intervalo-s}
\end{equation}
The inverse condition for the fixed reader is exactly
\begin{equation}
 s(v)=s_\alpha
 \quad\Longleftrightarrow\quad
 w(v)=\frac{73-s_\alpha}{24}
 \simeq0.00110309613175.
 \label{x_reciprocidad:lib02:criterio-alpha}
\end{equation}

None of the six states in \cref{x_reciprocidad:lib02:seis-estados} satisfies this condition. The second memory has \(s=73\), whereas the other five have \(s<72\). This last inequality is verified exactly as \(24b_v>a_v\). In the three rational rows, multiplying integers suffices; for \(u\) and \(q_K\), the differences are
\[
 24b_u-a_u=\frac{66073183}{20}+\frac{5235904}{5}\sqrt5>0,
\]
\[
 24b_{q_K}-a_{q_K}
 =\frac{45259241}{15}-\frac{5235904}{5}\sqrt5>0.
\]
For the second, \(\sqrt5<5/2\) already gives a positive lower bound. The exclusion concerns this reader and these states; it does not declare a global absence of relations with alpha.

Since \(0<s_\alpha<73\), the radius
\[
 r_\alpha=\frac{4\sqrt{s_\alpha}}{81}<\frac{4\sqrt{73}}{81}
\]
is a valid conservative guarantee, but it is not the optimal radius of the canonical reader. Any number between zero and 73 would likewise give a smaller guarantee. If \(a^2+a+1=s_\alpha\) were imposed by changing \(a\), then \(\sqrt{2a}\) and \((a+1)^2\) would also change; one cannot replace only 73 and retain 4 and 81 as though the operator had not changed.

\begin{theorem}[Exact weight iteration and reading recurrence]
\label{x_reciprocidad:lib02:iteraciones}
Let \(v_n=Q_{-,a}^nv_0\), \(a>0\), and \(b=(a-1)^2\). When \(v_n\ne0\),
\begin{equation}
 w_n=\frac{b^nw_0}{b^nw_0+p(a)^n(1-w_0)}.
 \label{x_reciprocidad:lib02:pesos-iterados}
\end{equation}
For \(a\ne1\) and \(0<w_0<1\),
\[
 \frac{w_n}{1-w_n}
 =\left(\frac b{p(a)}\right)^n\frac{w_0}{1-w_0},
 \quad w_n\downarrow0,\quad s_n\uparrow p(a),
\]
and
\begin{equation}
 s_{n+1}=p(a)+b-\frac{bp(a)}{s_n},\qquad
 s_{n+1}-s_n=\frac{(s_n-b)(p(a)-s_n)}{s_n}>0.
 \label{x_reciprocidad:lib02:riccati}
\end{equation}
\end{theorem}
\begin{proof}
The operator and its adjoints preserve the two orthogonal sectors. By \eqref{x_reciprocidad:lib02:gram-Qmenos}, the squared norm in the fixed sector is multiplied by \(b\) at each step, and that in the complementary sector by \(p(a)\). This gives \eqref{x_reciprocidad:lib02:pesos-iterados} and the weight ratio. Since \(p(a)-b=3a>0\), monotonicity and the limit are strict in the mixed case. Finally, \(s_n=\|v_{n+1}\|^2/\|v_n\|^2\); the sequence of norms is a combination of \(b^n\) and \(p(a)^n\), satisfying \(a_{n+2}=(p(a)+b)a_{n+1}-bp(a)a_n\). Division by \(a_{n+1}\) gives the recurrence, and factoring its difference gives \eqref{x_reciprocidad:lib02:riccati}.
\end{proof}

If \(a=1\), the fixed sector vanishes at the first step. When the
complement is nonzero, the subsequent reading is exactly three;
a purely fixed state is transformed into zero and its quotient ceases to
be defined. For \(a\ne1\), the inverse iteration satisfies
\(s_{n-1}=bp(a)/(p(a)+b-s_n)\) and decreases the mixed reading.
Reversing the orientation \(U\mapsto U^{-1}\), by contrast, leaves
the sequence of norms of the forward iterations unchanged.
Powers of \(Q_{+,a}\) multiply the weight ratio by
\(((a+2)^2/p(a))^n>1\) and decrease \(s_{-,a}\) in mixed
states. These are actual filters of the state, not chart changes or a
TPK temporal trajectory identified by definition.

\begin{proposition}[Complete test of all integer iterations]
\label{x_reciprocidad:lib02:cruces-enteros}
For \(a=8\), no integer iteration \(Q_-^nv\), \(n\geq0\), with \(v=K,k,m_0\), satisfies \(s=s_\alpha\). The only crossings of the target interval occur between indices \(21,22\), \(14,15\), and \(15,16\), respectively. For \(v=m_1\), the reading remains exactly 73 in every iteration.
\end{proposition}
\begin{proof}
Define the rational endpoints of the target weight:
\[
 w_L=\frac{73-h}{24}
 =\frac{2647430716199002714894527619337}
 {2400000000000000000000000000000000},
\]
\[
 w_H=\frac{73-\ell}{24}
 =\frac{441238452699833785815754603223}
 {400000000000000000000000000000000}.
\]
For \(w_0=A/(A+B)\), the exact formula is \(w_n=A49^n/(A49^n+B73^n)\). The data and crossings are
\[
\begin{array}{c|r|r|c}
 v&A&B&N\\\hline
 K&10684363&2069408&21\\
 k&3512283&8277632&14\\
 m_0&1052377&1837634&15
\end{array}
\]
and satisfy, in each row,
\[
 \frac{A49^N}{A49^N+B73^N}>w_H>w_L>
 \frac{A49^{N+1}}{A49^{N+1}+B73^{N+1}}.
\]
These are integer inequalities: for example, if \(w_H=p_H/q_H\), the first is equivalent to \((q_H-p_H)A49^N>p_HB73^N\), and the last is verified analogously with \(w_L\). Substitution of the three rows and the displayed integers proves the strict crossings without rounding alpha. Strict monotonicity in \cref{x_reciprocidad:lib02:iteraciones} excludes all preceding and subsequent indices; between \(N\) and \(N+1\) there is no other integer. For \(m_1\), \(w_0=0\) remains zero by \eqref{x_reciprocidad:lib02:pesos-iterados}.
\end{proof}

For numerical reference, the readings on either side of the crossings are
\[
\begin{array}{c|c|c}
 v&s_N&s_{N+1}\\\hline
 K&72.97136264842803&72.98077012438719\\
 k&72.96167997982852&72.97426483341773\\
 m_0&72.96527892867391&72.97668298511348
\end{array}
\]
The proof covers all \(n\geq0\), not only a finite sample.
The inverse orientation gives the same crossings; positive powers of
\(Q_+\) and inverse powers of \(Q_-\) decrease the
initial readings, already less than 72, and likewise do not reach the target
in these three families. The mixed products \(Q_-^nQ_+^m\)
constitute another two-index class, with weight ratio
\((49/73)^n(100/73)^m w_0/(1-w_0)\); they are not excluded by
the preceding one-parameter proof, nor are indices chosen here to fit
alpha.

\section{Angular reciprocity and the existing regional correction}
\label{x_reciprocidad:lib02:correccion-regional}

The chart changes of \cref{x_reciprocidad:lib02:inversion-carta} include the
reciprocity and elliptic deformation already constructed. In a plane
realized by \(M_\eta\), the transported form is
\(M_\eta^{-T}M_\eta^{-1}\); similarity does not change the
eigenvalues of the cycle. Retaining the previous Euclidean product after
changing the chart would define another observable. The involution
\(q_+\leftrightarrow q_-\) preserves \(A\), reverses \(C^*\)
and sends \(\eta_{\rm el}\) to \(-\eta_{\rm el}\). The polar
relation of the ellipse
\(\tan\phi=e^{-\eta_{\rm el}}\tan\theta\) changes the
parametrization of the curve, not the spectral phase of a similar
transport. It does not add a small phase to \(S^4\).

Incidence conjugation has a different type. Its block law preserves
\[
 z_m^\dagger=z_{\rho(m)}
 +100\mathbf1_{\{c_{\rho(m)}\ne0\}}-20c_{\rho(m)},
 \qquad c_m=[C_m]_{10},
\]
with the boundary terms before decimal reduction. In general, it is not an isometry of the block vector. Applying it to the register requires its connection to the incidence register and its endpoints; it is not replaced by an invented angular correction.

\begin{proposition}[Regional correction of the counter-angle]
\label{x_reciprocidad:lib02:retorno-regional}
The antecedent action and angle coordinates satisfy
\[
 A_{\rm deg}=1000\alpha_A,\qquad
 C^*_{\rm ret,deg}=2(\eta_{\rm ret}+\alpha_A),\qquad
 C^*_{\rm pre,deg}=2(H_5+\alpha_A),
\]
\[
 \eta_{\rm ret}=H_5-\frac{90}{\pi}\mathcal C_\pi(\alpha_A),
 \quad
 \mathcal C_\pi(\alpha_A)=169\alpha_A^6+
 \frac{D_A\alpha_A^7}{1-D_A\alpha_A},\qquad
 D_A=\exp(-100\pi\alpha_A/9).
\]
The difference in radians is exactly
\[
 \theta_{C,\rm pre}-\theta_{C,\rm ret}=\mathcal C_\pi(\alpha_A).
\]
\end{proposition}
\begin{proof}
The difference in degrees is \(2(H_5-\eta_{\rm ret})=180\mathcal C_\pi/\pi\). Multiplication once by \(\pi/180\) gives the result.
\end{proof}

Here \(\pi\) and \(\alpha_A\) are the coordinates already generated; \(H_5\) and the remaining readers retain their antecedent definitions, not a recalibration. The subsequent evaluations are
\[
\begin{array}{c|c}
 \text{coordinate}&\text{radians, approximately}\\\hline
 \theta_A&0.127362829013\\
 \theta_C&0.0370662264658\\
 \theta_C/6\text{ in each antipodal pair}&0.00617770441097\\
 \mathcal C_\pi&2.55207421805\cdot10^{-11}
\end{array}
\]
The channels \(q_\pm\) are positive real attenuations; their exponent does not automatically become a unitary phase through a factor \(i\). No additional divisors are introduced to approximate the scales of the following test.

\section{Angular inverse problem in the dodecaphasic register}
\label{x_reciprocidad:lib02:ensayo-angular}

For a general target \(s\in[49,81]\), the principal orientation of the pure sector is obtained from
\[
 |8-e^{i\theta}|^2=65-16\cos\theta,
 \qquad \theta(s)=\arccos\frac{65-s}{16}\in[0,\pi].
\]
In particular, the subsequent datum \(s_\alpha\) determines
\[
 \theta_\alpha=\theta(s_\alpha),\qquad
 \delta_\alpha=\frac{2\pi}{3}-\theta_\alpha.
\]
This formula solves an inverse problem using an output already produced; it does not claim to produce alpha anew.

\begin{lemma}[Angular structure of the ternary complement]
\label{x_reciprocidad:lib02:J}
The operator \(J=(U-U^2)/\sqrt3\) satisfies
\[
 J^*=-J,\quad J^2=-P,\quad JP_0=P_0J=0,\quad JP=PJ=J,
 \qquad U=P_0-\frac12P+\frac{\sqrt3}{2}J.
\]
\end{lemma}
\begin{proof}
The adjoint exchanges \(U,U^2\). Moreover, \((U-U^2)^2=U+U^2-2I=-3P\), and \((I+U+U^2)(U-U^2)=0\). The remaining relations follow from \(P=I-P_0\); adding \((U+U^2)/2\) and \((U-U^2)/2\) gives the final expression.
\end{proof}

\begin{theorem}[Angular family and preserved properties]
\label{x_reciprocidad:lib02:Utheta}
For \(\theta\in\mathbb R\), define
\[
 U_\theta=P_0+\cos\theta\,P+\sin\theta\,J,\qquad
 Q_{-,\theta}=8I-U_\theta.
\]
One has
\[
 U_\theta^*U_\theta=I,\qquad
 U_\theta U_\varphi=U_{\theta+\varphi},\qquad U_{2\pi/3}=U,
\]
\begin{equation}
 Q_{-,\theta}^*Q_{-,\theta}
 =49P_0+(65-16\cos\theta)P.
 \label{x_reciprocidad:lib02:gram-angular}
\end{equation}
The readers \(T_\theta=(8I+U_\theta)/9\), \(D_\theta=\sqrt8(U_\theta-I)/9\) preserve \(T_\theta^*T_\theta+D_\theta^*D_\theta=I\). By contrast,
\[
 U_\theta^3=I\ \Longleftrightarrow\ 3\theta\in2\pi\mathbb Z,
 \qquad
 \det Q_{-,\theta}=7^4(65-16\cos\theta)^4.
\]
\end{theorem}
\begin{proof}
The relations in \cref{x_reciprocidad:lib02:J} cancel the terms between sectors. On the complement, \(J^2=-I\), so that \(U_\theta^*U_\theta=P_0+(\cos^2\theta+\sin^2\theta)P=I\). Multiplying two expressions and using the addition formulas gives the composition law. The expression for \(U\) in the lemma gives the ternary value. Since \(U_\theta+U_\theta^*=2P_0+2\cos\theta P\), expanding the Gram matrix gives \eqref{x_reciprocidad:lib02:gram-angular}. The expansions
\[
 81T_\theta^*T_\theta=65I+8(U_\theta+U_\theta^*),\qquad
 81D_\theta^*D_\theta=16I-8(U_\theta+U_\theta^*)
\]
prove the balance. The composition law gives \(U_\theta^3=P_0+\cos(3\theta)P+\sin(3\theta)J\), which is the identity exactly under the indicated condition.

For the determinant, \(Q_{-,\theta}\) acts by 7 on the four fixed dimensions. On the complement, a unit vector \(e\) and \(Je\) form an orthonormal pair; their orthogonal complement is stable under \(J\). Repeating this yields four planes, on each of which the determinant is \((8-\cos\theta)^2+\sin^2\theta=65-16\cos\theta\). Multiplying the contributions proves the formula.
\end{proof}

\begin{corollary}[Inverse solution and separation of corrections]
\label{x_reciprocidad:lib02:fase-alpha}
On \(\operatorname{im}P\), the reader \(Q_{-,\theta_\alpha}\) has quadratic ratio exactly \(s_\alpha\). The orientations \(\pm\theta_\alpha\) give the same ratio, and
\[
 73-s_\alpha
 =8(1-\cos\delta_\alpha+\sqrt3\sin\delta_\alpha).
\]
We have \(U_{\theta_\alpha}^3\ne I\), and \(\delta_\alpha\ne\mathcal C_\pi(\alpha_A)\).
\end{corollary}
\begin{proof}
The definition of the phase solves for the complementary coefficient in \eqref{x_reciprocidad:lib02:gram-angular}; cosine is even. Substituting \(\theta_\alpha=2\pi/3-\delta_\alpha\) gives the stable identity for the difference. Since \(72<s_\alpha<73\), \(\pi/2<\theta_\alpha<2\pi/3\) and \(0<3\delta_\alpha<\pi/2\); hence \(3\theta_\alpha\) is not a multiple of \(2\pi\).

The separation of corrections can be proved without relying on approximate figures. The interval gives \(73-s_\alpha>0.026\). If \(\delta_\alpha\leq10^{-3}\), the inequalities \(1-\cos t\leq t^2/2\), \(\sin t\leq t\), and \(\sqrt3<2\) would give
\[
 73-s_\alpha\leq8(10^{-6}/2+2\cdot10^{-3})=0.016004,
\]
a contradiction. Thus \(\delta_\alpha>10^{-3}\). On the other hand, \(0<\alpha_A<10^{-2}\), \(0<D_A<1\) imply
\[
 0<\mathcal C_\pi(\alpha_A)
 <169\cdot10^{-12}+\frac{10^{-14}}{1-10^{-2}}
 <170\cdot10^{-12}<10^{-7}.
\]
The two quantities are distinct.
\end{proof}

The diagnostic evaluation gives
\[
 \delta_\alpha\simeq0.00190956706826\ {\rm rad}
 \simeq0.109410133709^\circ,\qquad
 \theta_\alpha\simeq119.890589866291^\circ.
\]
The regional correction is approximately \(7.48\cdot10^7\) times
smaller. The inverse family preserves orthogonality and balance, but loses
order three at that phase and changes the original determinant. Its real
chart does not by itself provide an integer matrix or the lattice quotient of
index 73. Solving the prescribed datum does not prove
\(73=10^4\alpha_A\).

\begin{proposition}[Reachability with the state preserved]
\label{x_reciprocidad:lib02:alcanzabilidad}
For \(v\ne0\), with \(w=w(v)\) fixed,
\[
 s_\theta(v)=49w+(65-16\cos\theta)(1-w).
\]
If \(w<1\), its image is exactly \([49,81-32w]\). Hence a phase attaining \(s_\alpha\) exists if and only if
\[
 w\leq\frac{81-s_\alpha}{32}\simeq0.250827322099.
\]
When it exists, its principal orientation is
\[
 \theta_v=\arccos\left[
 \frac1{16}\left(65-\frac{s_\alpha-49w}{1-w}\right)\right].
\]
If \(w=1\), the reading is always 49 and does not reach the target.
\end{proposition}
\begin{proof}
Evaluating \eqref{x_reciprocidad:lib02:gram-angular} and dividing by \(\|v\|^2\) gives the formula. Cosine ranges over \([-1,1]\), so the complementary coefficient ranges over \([49,81]\). Its combination with the fixed weight gives the entire interval, without gaps. Solving the equality and checking the domain of arccosine gives the condition and the phase. The case \(w=1\) is evaluated without dividing by \(1-w\).
\end{proof}

\begin{corollary}[Four impossibilities and two possible phases]
\label{x_reciprocidad:lib02:seis-angulares}
Keeping the state and the fixed sector, the angular family does not attain
\(s_\alpha\) on \(K,k,m_0,u\). It does attain it on
\(m_1\) and \(q_K\).
\end{corollary}
\begin{proof}
From the exact table \eqref{x_reciprocidad:lib02:tabla-normas}, we obtain
\[
\begin{array}{c|c}
 v&32b_v-9a_v\\\hline
 K&227115677/3\\
 k&2094607/4\\
 m_0&122655440/243\\
 u&37201759/20+(7859008/15)\sqrt5
\end{array}
\]
All entries are positive. Thus \(w>9/32\), and the maximum \(81-32w\) is less than \(72<s_\alpha\). For \(m_1\), \(b_{m_1}=0\) and \(a_{m_1}>0\), so the interval is \([49,81]\). For \(q_K\),
\[
 a_{q_K}-4b_{q_K}=\frac{588839+1195712\sqrt5}{15}>0.
\]
Hence \(w(q_K)<1/4\) and its maximum is greater than \(73>s_\alpha\). In both cases, the target belongs to the reachable interval.
\end{proof}

The evaluations of the maxima and, where they exist, the phases are
\[
\begin{array}{c|c|c}
 v&\max_\theta s_\theta(v)&\theta_v\\\hline
 K&54.1922726227&\text{does not exist}\\
 k&71.4670172771&\text{does not exist}\\
 m_0&69.3474270513&\text{does not exist}\\
 u&66.8296175021&\text{does not exist}\\
 m_1&81&119.8905898663^\circ\\
 q_K&77.3923674871&133.5296853010^\circ
\end{array}
\]
The phase of the pure sector cannot be used indiscriminately for a mixed
state: its reading would be
\(49w+s_\alpha(1-w)=s_\alpha-(s_\alpha-49)w\).
The original readers of \(V_2\) are polynomials in \(S\), such as
\(P_0,P,J,U_\theta\). They therefore also intertwine the angular
family, and \(s_\theta(V_2v)=s_\theta(v)\). Archiving the
complements preserves the joint reading, but does not remove the
reachability restrictions.

\section{Composition of readers and scope of inverse problems}
\label{x_reciprocidad:lib02:alcance}

The factor \(73\) has been determined by three compatible
and distinct operations: the Gram matrix of two nonadic readers, the
similarity of the ternary augmentation module and the integer determinants of
the two conjugate charts. Bilateral balance preserves the complete
state and admits explicit inversion; complete memory preserves the
transported readings, while selecting a channel can change
their weights. The correction radii are uniform over all integer
registers in the algebraic domain and throughout the family \(a>0\), with
cut and parity proofs that do not depend on sampling.

The inverse test adds verifiable conditions, not a hidden
calibration. Orientation inversion and metric covariance do not change
the reading. The specified integer iterations are resolved for
all their indices by an exact crossing and monotonicity. Angular variation
produces a family that realizes the target in the allowed states,
but the datum \(s_\alpha\) selects its phase and the original ternary
index is lost. Its identification with a prior realization of
\(\Gamma_9\) would require a typed composition on the enriched
state; this test neither provides such an identification nor declares its
absence throughout the corpus.

The new formalizations assemble antecedent operators and questions; no claim of worldwide priority is made for encoding, the use of cyclic vectors, or polar decomposition. The decimal figures presented are subsequent evaluations of exact formulas. The conditions of nonvanishing, charge, metric, lattice, and normalization are part of the result, and their counterexamples remain alongside the proofs. The function of the holographic arithmetic-geometric coupling constant is thereby preserved: a produced object that articulates readers, memory, and incidence, without turning every numerical coincidence into a new physical law.
